% Source recovery: noncommutative central cubic pair, all Artin bases.
% Main mathematical input: original 2026-11-01 paper V, sections 12--13 and 21.
% This section is included in papers/nc_cubic_threefold_moduli.tex.
\section{中心三次対の全Artin形式局所構造と二次包絡}
\label{sec:moduli-artin}
本節は、滑らかな可換三次三様体の近傍において、
有限係数の17次元族と異なる\emph{完備形式局所モジュライ}を扱う。
以前の中心三次対原稿の第12--13節から定理・証明の有効部分を
再構成したものである。主張は生成空間に標識を入れた
通常の集合値Artin変形関手に関するものであり、
derived stackの同型でも非可換領域全体の大域分類でもない。
外部の同次formality、Ore局所化およびAbel圏変形の
適用条件について、原稿に記載した証明と今回の独立監査を区別する。

以下$k=\C$、$V=k^5$、$S=\Sym V^*$、
$F\in S_3$を滑らかな三次三様体
$X=V_+(F)\subset\PP(V)$の定義式とする。
$(R,\mathfrak m)$は剰余体$k$の局所Artin代数である。

\begin{lemma}[平坦な正規三次商は自動的に中心商]
\label{lem:artin-normal-central}
$B_R$を$R$上各次数有限自由な連結非負次数付き代数とし、
一次で生成され、特殊ファイバー$B_R/\mathfrak mB_R=S$とする。
$w\in(B_R)_3$が$F$に特殊化するなら$w$は左右正則であり、
以下の条件は同値である：
\[
 B_R/(w)\text{ が各次数で }R\text{-平坦},\qquad
 wB_R=B_Rw,\qquad
 w\in Z(B_R).
\]
成立するとき商のHilbert級数は$(1-t^3)/(1-t)^5$である。
\end{lemma}
\begin{proof}
$w$の左右乗法の各次数成分は、$\mathfrak m$で還元すると
$F:S_{d-3}\hookrightarrow S_d$である。
$F\ne0$と$S$が整域であることから単射となり、
有限自由行列の最大rank小行列式が単元であるため、
元の乗法も分裂単射になる。

$J=(w)$について商が平坦なら
$0\to J_d\to(B_R)_d\to(B_R/J)_d\to0$は
有限自由$R$加群の分裂完全列である。
$J_d/\mathfrak mJ_d=FS_{d-3}$なので、
左右乗法$(B_R)_{d-3}\to J_d$はNakayamaと
先の単射性により同型である。従って$J=wB_R=B_Rw$。
正規性から$wb=\sigma(b)w$で次数付き自己同型$\sigma$が得られる。
特殊ファイバーでは$\bar\sigma=\mathrm{id}$であり、
$b=w$と右消去によって$\sigma(w)=w$となる。
$\sigma\equiv\mathrm{id}\pmod{\mathfrak m^r}$とし、
$\mathfrak m^r/\mathfrak m^{r+1}=I$とおくと、
次の差分は次数零の導分
\[
 D:S\longrightarrow I\otimes_k S,\qquad
 D(F)=0
\]
である。$D=\sum_i\ell_i\partial_{x_i}$、
$\ell_i\in I\otimes S_1$だから
$\sum_i\ell_iF_i=0$。滑らかさから
$F_0,\ldots,F_4$は二次正則列をなし、
Koszul複体に一次係数のsyzygyが存在しない。
よって$D=0$となり、$\mathfrak m$の冪零性から
$\sigma=\mathrm{id}$を得る。
中心性から商平坦性へは、
各次数の分裂単射$w:(B_R)_{d-3}\to(B_R)_d$を使えばよい。
\end{proof}

\begin{corollary}[標識付き対の自己同型]
\label{cor:artin-scalar}
上の中心正規対$(B_R,Rw)$の次数付き自己同型が
特殊ファイバーで恒等なら、次数一の生成元への作用は
すべて同じ単元$\lambda\in1+\mathfrak m$によるスカラー倍である。
従って射影標識を入れた対には非自明な無限小自己同型はない。
\end{corollary}
\begin{proof}
$\phi(w)=cw$、$c\in1+\mathfrak m$と書き、
$c$の一意の三乗根による次数スカラー倍を合成すると
$w$を固定する自己同型に帰着する。
その$\mathfrak m$進差分は前補題の$D(F)=0$と同じ
一次Koszul syzygyで消える。
\end{proof}

\subsection{相対Casimir条件と五つのPl\"ucker方程式}
体積形式$\Omega=dx_0\wedge\cdots\wedge dx_4$を固定する。
$F_R\in R\otimes S_3$を$F$の持上げとし、
$a=\sum_{i<j}a_{ij}\,dx_i\wedge dx_j$を
係数$a_{ij}\in\mathfrak m$の定数二形式とする。
\[
 \iota_{\Pi(F_R,a)}\Omega=dF_R\wedge a
\]
で二次bivectorを定める。
$\{i,j,k,l\}=\{0,\ldots,4\}\setminus\{r\}$、
$i<j<k<l$に対し
\[
 q_r(a)=(-1)^r
 (a_{ij}a_{kl}-a_{ik}a_{jl}+a_{il}a_{jk}),\qquad
 g=\sum_{r=0}^4q_r(a)\,\partial_rF_R
\]
とおく。

\begin{lemma}[全Artin上でのCasimir核]
\label{lem:artin-casimir}
二次bivector$\Pi$が$\iota_{dF_R}\Pi=0$を満たすことと、
一意な$a\in R\otimes\wedge^2V^*$を用いて
$\Pi=\Pi(F_R,a)$と表せることは同値である。
この同型は任意の$R$の基底変換と可換であり、
$\Pi$は発散零である。
\end{lemma}
\begin{proof}
$F$の五つの偏微分は二次正則列だから、そのKoszul複体の
\[
 \wedge^3V\otimes R
 \xrightarrow{\iota_{dF_R}}
 S_{R,2}\otimes\wedge^2V
 \xrightarrow{\iota_{dF_R}}
 S_{R,4}\otimes V
\]
の該当次数は、特殊ファイバーで左端の単射と中央での完全性を持つ。
有限自由加群の複体なので、分裂小行列式の単元性を
Artin基底へ持ち上げて同じ完全性を得る。
$\wedge^3V\simeq\wedge^2V^*$により表示と一意性が従う。
$d(dF_R\wedge a)=0$から発散零である。
\end{proof}

\begin{proposition}[JacobiイデアルとPl\"uckerイデアルの一致]
\label{prop:artin-pfaffian}
$\Pi=\Pi(F_R,a)$に対し、Jacobi方程式の係数イデアルは
\[
 (q_0(a),\ldots,q_4(a))
\]
と一致する。特に集合論的なrank分類だけでなく、
非被約な局所Artin基底上での方程式を同定する。
\end{proposition}
\begin{proof}
$J=\frac12[\Pi,\Pi]$と書くと、五変数の外積計算から
\[
 \iota_{\Pi\wedge\Pi}\Omega=2g\,dF_R,\qquad
 \iota_J\Omega=dg\wedge dF_R
\]
が得られる。これは多項式恒等式である。
符号は$ a=dx_0\wedge dx_1+dx_2\wedge dx_3$
の代入で直接確認できる。
線形写像
\[
 L_F:k^5\longrightarrow\Omega^2_{S,3},\qquad
 (c_r)\longmapsto d\bigl(\textstyle\sum_rc_rF_r\bigr)\wedge dF
\]
は単射である。核の元について、
$F_i$のKoszul完全性から
$d(\sum c_rF_r)$は$dF$の多項式倍となるが、
係数の次数は負なので零である。
Euler式と定数係数のKoszul syzygyの消滅により$c_r=0$。
従って$L_F$の非零$5$次小行列式は
$R$上でも単元であり、$L_{F_R}$は分裂単射となる。
上の$J$の表示の係数から$q_r(a)$を回収でき、
両方のイデアルが一致する。
\end{proof}

\begin{theorem}[中心三次対の全Artin形式局所環：入力を明示した結論]
\label{thm:artin-central-pair}
局所Artin基底上の次数別平坦な二次多項式環の標識付き変形と
同次Poisson変形との間に、固定したGL$_5$同変の
Kontsevich形式量子化を適用する。
中心三次直線が$kF$へ特殊化する対の古典的変形関手は
\[
 \widehat{\mathcal I}_{F}\simeq
 \operatorname{Spf}\frac{
  k[[s_1,\ldots,s_{34},a_{01},\ldots,a_{34}]]
 }{(q_0(a),\ldots,q_4(a))}
\]
で表される。形式局所次元は$41$、接空間次元は$44$。
特殊ファイバーで恒等に還元する射影座標変更を割ると、
\[
 \operatorname{Spf}\frac{
 k[[t_1,\ldots,t_{10},a_{01},\ldots,a_{34}]]
 }{(q_0(a),\ldots,q_4(a))}
\]
となり、次元$17$、接空間次元$20$である。
有限群$\Aut(X)$はこの記述ではまだ割らない。
\end{theorem}
\begin{proof}
同次形式量子化の星積を$B_\Pi$とする。
$\Pi$に対するtwisted Taylor mapの次数三部分
$D^0_\Pi:S_{R,3}\to(B_\Pi)_3$は
還元が恒等の有限自由加群の同型である。
また次数四のベクトル場$Y$の評価
\[
 D^1_\Pi(Y)\longmapsto
 \bigl(D^1_\Pi(Y)(x_0),\ldots,D^1_\Pi(Y)(x_4)\bigr)
\]
は同じく還元が恒等の同型なので、
twisted chain identityにより
\[
 D^0_\Pi(F_R)\in Z(B_\Pi)
 \quad\Longleftrightarrow\quad [\Pi,F_R]=0
\]
となる。ここで中心の形成と非平坦基底変換の交換性は
仮定せず、有限自由行列の可逆性によって逆向きを証明する。
このformalityの恒等式は
\href{https://arxiv.org/abs/math/0012228}{Cattaneo--Felder--Tomassini}
のtwisted formulaを入力としている。
補題\ref{lem:artin-normal-central}と
命題\ref{prop:artin-pfaffian}により、
$F_R$の射影変形の34変数と$a$の10変数、
五つのPl\"ucker方程式が得られる。
後者は$\mathrm{Gr}(2,5)$の非零アフィン錐の閉包で
次元$7$、原点の接空間次元$10$なので
$34+7=41$、$34+10=44$。
$X$が滑らかで$H^0(X,T_X)=0$だから
$\PGL_5$の軌道方向24は形式的に横断sliceから分離できる。
$34-24=10$を得て後半が従う。
なお特異な形式基底でのCalabi--Yau性は相対的であり、
絶対大域次元が五であるとは主張しない。
\end{proof}

\begin{remark}[一次方向はすべて二次へ延長するわけではない]
\label{rem:artin-obstruction}
$a=\varepsilon(dx_0\wedge dx_1+dx_2\wedge dx_3)$を
$k[\varepsilon]/(\varepsilon^2)$上で考えると
$q_r(a)=0$だが、
$k[\varepsilon]/(\varepsilon^3)$上では$q_4(a)=\varepsilon^2$。
二次の補正$b$を$a+\varepsilon^2b$に加えてもこの項は消えない。
従って接空間20次元と局所次元17の差は実際の二次障害を表し、
全Artin比較を無障害性と混同してはいけない。
\end{remark}

\subsection{Ore局所化、偏極、および斉次二次包絡}
次数別平坦な斉次変形$A_R$を
$S/(F)$へ特殊化するものとする。
一次元の持上げ$\widetilde x_i$を選び、
標準開集合$U_i=D_+(x_i)\cap X$上で
\[
 \mathscr A_R(U_i)=
 \bigl(A_R[\widetilde x_i^{-1}]\bigr)_0
\]
を取る。$[A_R,A_R]\subset\mathfrak m A_R$かつ
$\mathfrak m$冪零なので、有限の交換子展開により
$\{\widetilde x_i^n\}$は左右Ore集合となる。
その元の特殊化は$S/(F)$の整域の非零因子だから、
各局所化は$R$平坦である。
重なりは二つの一次元をともに可逆化した次数零部分で
記述され、三重重なりのcocycleも局所化の普遍性に従う。
この操作で$X$上に平坦な結合$R$代数層
$\mathscr A_R$が得られる。

\begin{proposition}[偏極とsection algebraの回復]
\label{prop:artin-ore-section}
上のOre局所化層の次数$n$部分を貼り合わせた
$\mathscr E_n$は可逆両側$\mathscr A_R$加群であり、
\[
 \mathscr E_m\otimes_{\mathscr A_R}\mathscr E_n
 \simeq\mathscr E_{m+n}.
\]
$n\ge0$について
\[
 H^i(X,\mathscr E_n)=0\ (i>0),\qquad
 (A_R)_n\xrightarrow{\sim}H^0(X,\mathscr E_n)
\]
が成り立ち、積とも両立し任意のArtin基底変換に可換である。
\end{proposition}
\begin{proof}
局所化の一次元は可逆なので強い$\mathbb Z$次数付けと
両側可逆性が従う。有限標準被覆の\v{C}ech複体を取り、
$R$の小拡大を順に帰納すると、
特殊ファイバーの$H^i(X,\OO_X(n))=0$ ($i>0,n\ge0$)
から高次消滅と$H^0$の有限自由性が得られる。
特殊化で同型な$(A_R)_n\to H^0(X,\mathscr E_n)$は
Nakayamaで同型となる。強次数付けにより積との両立も従う。
\end{proof}

\begin{lemma}[可逆両側偏極の同型類の一意性]
\label{lem:artin-polarization}
$\mathscr A_R$上の二つの平坦な可逆両側加群
$\mathscr M,\mathscr N$が同じ線束$L$に特殊化し、
$H^1(X,\OO_X)=H^0(X,T_X)=0$なら、
$\mathscr M\simeq\mathscr N$である。
$H^0(X,\OO_X)=k$なら各々の自己同型は$R^\times$である。
\end{lemma}
\begin{proof}
小拡大$R'\twoheadrightarrow R$の核$I$で
$\mathfrak mI=0$とする。
特殊ファイバーの局所基底の持上げにより
$\mathscr M,\mathscr N$は右加群として局所自由階数一。
同型の右加群としての持上げの障害は
$I\otimes H^1(X,\OO_X)$に属し消える。
二つの左作用の差は導分に対応し
$I\otimes H^0(X,T_X)=0$なので持上げは両側同型。
自己同型群は
$H^0(X,\OO_X)=k$と小拡大帰納から$R^\times$となる。
\end{proof}

\begin{theorem}[斉次環の全Artin変形の二次cover：比較定理を入力]
\label{thm:artin-quadratic-cover}
上の中心三次対の形式関手と、
$X$上の平坦結合代数層の標識付き古典的Artin変形関手の間の
Ore局所化による自然変換$\Theta$を考える。
YekutieliのArtin変形量子化の自然性と
De~Deken--Lowenの$\mathbb Z$-生成列の変形比較を
この設定に適用すると、$\Theta$は自然同型となる。
従って任意の次数別平坦な斉次変形$A_R$は
\[
 V_R=(A_R)_1,\quad
 K_{2,R}=\ker(V_R^{\otimes2}\to(A_R)_2),\quad
 B_R^{\mathrm{quad}}=T_R(V_R)/(K_{2,R})
\]
を用いて一意の中心三次対から回復される。
特に$B_R^{\mathrm{quad}}$は次数別$R$平坦で、
\[
 L_R=\ker\bigl((B_R^{\mathrm{quad}})_3\to(A_R)_3\bigr)
 \subset Z(B_R^{\mathrm{quad}})_3
\]
は階数一の直和因子であり、
$A_R\simeq B_R^{\mathrm{quad}}/(L_R)$。
この構成はArtin基底変換と可換である。
\end{theorem}
\begin{proof}[原研究稿の証明構造]
まず局所化による$\Theta$を上で実際に構成した。
中心対の形式関手は
定理\ref{thm:artin-central-pair}の完備局所環で表される。
構造代数層のArtin変形関手も、
\href{https://arxiv.org/abs/0905.0488}{Yekutieli, Cor.~11.2}
の量子化と可換三次の変形sliceから同じ環で表される。
接写像$d\Theta$は、可換方向では
$H^1(X,T_X)$上の恒等写像であり、
反対称交換子方向では
\[
 \wedge^2V^*\longrightarrow H^0(X,\wedge^2T_X),
 \qquad a\longmapsto\iota_Ea
\]
という10次元空間の同型である。
よって接写像は可逆な三角行列となる。
完備局所環の連続自己写像が余接空間で同型なら
各$\mathfrak m^r/\mathfrak m^{r+1}$上でも同型であり、
形式的Nakayama帰納から環自体が同型となる。
したがって$\Theta$は全Artin基底で自然同型。

一方、上の命題で$\mathscr A_R$の局所化次数一は
$\OO_X(1)$の両側可逆な持上げ$\mathscr E_1$となり、
全section algebraを回復する。
補題\ref{lem:artin-polarization}でこの持上げの同型類は
一意なので、任意の$A_R$と$\Theta$から来る中心商を
次数付き環として同一視できる。
中心商では三次まで新しい二次関係がないから
二次包絡は元の$B_R$に一致し、三次核は$L_R$。
この方法は未知の二次包絡の平坦性を先取りせず、
既に存在する中心対に同型化した後でそれを導く。
基底変換は有限自由の分裂核と局所化の普遍性で従う。

ここで使用した一次資料は
\href{https://arxiv.org/abs/1209.6243}{Yekutieli,
 \emph{Deformations of affine varieties and the Deligne crossed groupoid}}
のOre局所化、
\href{https://arxiv.org/abs/math/0511197}{Lowen,
 \emph{Algebroid prestacks and deformations of ringed spaces}}、
\href{https://arxiv.org/abs/1001.4265}{De~Deken--Lowen,
 \emph{Abelian and derived deformations in the presence of
 $\mathbb Z$-generating geometric helices}}である。
これらの定理の適用範囲の独立した一次資料監査は未完成である。
\end{proof}

\begin{remark}[明確に主張しない事柄]
\label{rem:artin-limits}
この全Artin同型は偏極付き\emph{古典的}変形関手の同定であり、
derived moduli stack、高次自己同型を保つgroupoid、
構造対象を忘れたKuznetsov成分の変形、
またはそこで構成される可換復元との一致を与えない。
無限小の中心対を得ることと、有限非可換パラメータで
滑らかなtailsを持つことは別問題である。
\end{remark}

\subsection{固定ambientからの忘却の四次元ファイバー}
\begin{proposition}[中心三次直線の射影束]
\label{prop:artin-fixed-ambient-fiber}
ある良い開集合$U$の各点$B$が二次Koszul AS正則CY五次元整域であり、
$\dim Z(B)_1=2$、$\dim Z(B)_3=5$かつ
中心の形成が基底変更と可換であると仮定する。
$E_1=Z(\mathscr B)_1$、$E_3=Z(\mathscr B)_3$とおけば
\[
 0\longrightarrow\Sym^3E_1\longrightarrow E_3
 \longrightarrow L\longrightarrow0,\qquad\rank L=1.
\]
中心三次直線のパラメータ空間は
$\PP_U(E_3)\to U$で各fiberは$\PP^4$。
$\PP(E_3)\setminus\PP(\Sym^3E_1)$は
$\Sym^3E_1\otimes L^{-1}$を作用群とするtorsorである。
同型類を取る場合、$\Gamma_B=\Aut_{\gr}(B)/k^\times$
が有限なら一般fiberの粗い双有理型は
$\PP Z(B)_3/\Gamma_B$となり、その次元は4である。
\end{proposition}
\begin{proof}
中心三次部のベクトル束の射影化で前半を得る。
$B$が整域なので中心三次の非零元は正則、
次数ごとの乗法の核は零で、商の各次数が局所自由となる。
$E_3\to L$に非零に写る直線はこの短完全列のsplittingに対応し、
splittingの差は$\Hom(L,\Sym^3E_1)$に属する。
同型類については、二次包絡と三次核の復元から
$B/(W)\simeq_{\gr}B/(W')$と
$[W'],[W]$が$\Gamma_B$軌道に属することが同値である。
Rosenlichtの有理商を$U$とincidenceへ適用し、
周囲$B$の軌道の上の空間が
$\PGL(V)\times_{\Gamma_B}\PP Z(B)_3$となることから
不変関数体が$k(\PP Z(B)_3)^{\Gamma_B}$に一致する。
有限群商は次元を保つ。
\end{proof}

\begin{remark}[四係数は標識を忘れても消えない]
特に固定した周囲に対し$Z(B)_3=kG\oplus\Sym^3Z(B)_1$なら、
$G+f(u,v)$の四係数が相対アフィンfiberを表す。
有限群$\Gamma_B$による同一視は生じ得るが、
四係数全体を座標変換で消せるわけではない。
原稿中の明示剛体周囲$B_*$についての
$\Aut_{\gr}(B_*)=k^\times$と正則開集合の主張は、
付属の有限整数証明書が現行Git treeに存在しないため、
本稿では独立に再検証した定理として採用しない。
\end{remark}
